\documentclass[12pt,letterpaper]{article}

% ── Encoding & fonts ──────────────────────────────────────────────────────────
\usepackage[T1]{fontenc}
\usepackage[utf8]{inputenc}
\usepackage{lmodern}

% ── Page layout ───────────────────────────────────────────────────────────────
\usepackage[margin=1in]{geometry}

% ── Graphics ──────────────────────────────────────────────────────────────────
\usepackage{graphicx}
\graphicspath{{./images/}}

% ── Tables ────────────────────────────────────────────────────────────────────
\usepackage{booktabs}
\usepackage{tabularx}
\usepackage{array}
\usepackage{longtable}

% ── Colours ───────────────────────────────────────────────────────────────────
\usepackage[dvipsnames]{xcolor}
\definecolor{awsorange}{RGB}{255,153,0}
\definecolor{gcpblue}{RGB}{66,133,244}
\definecolor{azureblue}{RGB}{0,120,212}
\definecolor{codebg}{RGB}{245,245,245}
\definecolor{coderule}{RGB}{200,200,200}

% ── Hyperlinks ────────────────────────────────────────────────────────────────
\usepackage[
  colorlinks=true,
  linkcolor=black,
  urlcolor=blue,
  citecolor=black,
  pdfborder={0 0 0}
]{hyperref}

% ── Code listings ─────────────────────────────────────────────────────────────
\usepackage{listings}
\lstdefinestyle{jsonStyle}{
  language=,
  backgroundcolor=\color{codebg},
  frame=single,
  rulecolor=\color{coderule},
  basicstyle=\ttfamily\small,
  breaklines=true,
  captionpos=b,
  numbers=none,
  showstringspaces=false,
  aboveskip=6pt,
  belowskip=6pt,
}
\lstdefinestyle{bashStyle}{
  language=bash,
  backgroundcolor=\color{codebg},
  frame=single,
  rulecolor=\color{coderule},
  basicstyle=\ttfamily\small,
  keywordstyle=\color{blue},
  commentstyle=\color{gray},
  breaklines=true,
  captionpos=b,
  numbers=none,
  showstringspaces=false,
  aboveskip=6pt,
  belowskip=6pt,
}

% ── Section styling ───────────────────────────────────────────────────────────
\usepackage{titlesec}
\titleformat{\section}{\large\bfseries}{}{0pt}{}[\titlerule]
\titleformat{\subsection}{\normalsize\bfseries}{}{0pt}{}
\titleformat{\subsubsection}{\normalsize\itshape}{}{0pt}{}

% ── Lists ─────────────────────────────────────────────────────────────────────
\usepackage{enumitem}
\setlist[itemize]{noitemsep, topsep=4pt}
\setlist[enumerate]{noitemsep, topsep=4pt}

% ── Misc ──────────────────────────────────────────────────────────────────────
\usepackage{microtype}
\usepackage{parskip}
\setlength{\parskip}{6pt}
\usepackage{fancyhdr}
\setlength{\headheight}{14pt}
\pagestyle{fancy}
\fancyhf{}
\lhead{\small CS 5525 --- Cloud Computing}
\rhead{\small Tony Nguyen}
\cfoot{\small\thepage}
\renewcommand{\headrulewidth}{0.4pt}

% ── Screenshot helper ─────────────────────────────────────────────────────────
% Usage: \screenshot{rel/path}{caption text}
\newcommand{\screenshot}[2]{%
  \begin{figure}[htbp]
    \centering
    \includegraphics[width=0.85\linewidth]{#1}
    \caption{#2}
  \end{figure}%
}

% ═════════════════════════════════════════════════════════════════════════════
\begin{document}

% ── Title page ────────────────────────────────────────────────────────────────
\begin{titlepage}
  \centering
  \vspace*{2cm}
  {\LARGE\bfseries Assignment 2 --- Multi-Cloud Web Deployment\par}
  \vspace{1cm}
  \rule{0.6\linewidth}{0.4pt}
  \vspace{1cm}

  {\large
    \textbf{Course:} CS 5525 --- Cloud Computing\\[6pt]
    \textbf{Student:} Tony Nguyen\\[6pt]
    \textbf{Date:} March 17, 2026\\[6pt]
    \textbf{University:} University of Missouri -- Kansas City
  }

  \vspace{2cm}
  \rule{0.6\linewidth}{0.4pt}
  \vspace{1cm}

  {\normalsize
    Deployment across \textbf{Amazon Web Services (AWS)},\\
    \textbf{Google Cloud Platform (GCP)}, and \textbf{Microsoft Azure}
  }
  \vfill
\end{titlepage}

% ── Table of Contents ─────────────────────────────────────────────────────────
\tableofcontents
\newpage

% =============================================================================
\section{Overview}
% =============================================================================

This assignment demonstrates deploying a personal web page across three major cloud
providers --- \textbf{Amazon Web Services (AWS)}, \textbf{Google Cloud Platform (GCP)},
and \textbf{Microsoft Azure} --- using each platform's preferred hosting method. The
goal is to compare Infrastructure-as-a-Service (IaaS) vs.\ Platform-as-a-Service (PaaS)
approaches, understand IAM and access control policies, and practice real-world cloud
deployment workflows.

% =============================================================================
\section{Live Deployments}
% =============================================================================

\begin{center}
\begin{tabularx}{\linewidth}{@{} l X l @{}}
\toprule
\textbf{Platform} & \textbf{URL} & \textbf{Hosting Method} \\
\midrule
\textbf{AWS} &
  \href{http://amzn-s3-cs5525-tonyn-bucket.s3-website-us-east-1.amazonaws.com}%
       {\small amzn-s3-cs5525-tonyn-bucket.s3-website-us-east-1.amazonaws.com} &
  S3 Static Website \\[4pt]
\textbf{GCP} &
  \href{http://34.66.74.246}{\small http://34.66.74.246} &
  Compute Engine + Apache2 \\[4pt]
\textbf{Azure} &
  \href{https://cs5525tonynstorage.z19.web.core.windows.net}%
       {\small cs5525tonynstorage.z19.web.core.windows.net} &
  Blob Storage Static Website \\
\bottomrule
\end{tabularx}
\end{center}

% =============================================================================
\section{Web Page}
% =============================================================================

A simple personal portfolio HTML page (\texttt{index.html}) was created and deployed
on all three platforms, featuring:

\begin{itemize}
  \item Full name: \textbf{Tony Nguyen}
  \item Title: \textit{Data Scientist \& Cloud Architect}
  \item About Me section
  \item Skills and project highlights
\end{itemize}

% =============================================================================
\section{Platform Deployment Details}
% =============================================================================

% ─────────────────────────────────────────────────────────────────────────────
\subsection{1.\ Amazon AWS (30\%)}
% ─────────────────────────────────────────────────────────────────────────────

\textbf{Method:} S3 Static Website Hosting \textit{(PaaS)}

\subsubsection{Steps Taken}

\begin{enumerate}
  \item Signed into the \textbf{AWS Management Console} at
        \texttt{console.aws.amazon.com}.
  \item Navigated to \textbf{S3} $\to$ Created a new bucket:
        \texttt{amzn-s3-cs5525-tonyn-bucket} (Region: us-east-1).
  \item Enabled \textbf{Static Website Hosting} under the bucket
        \textbf{Properties} tab.
  \item Set \texttt{index.html} as the index document.
  \item Disabled ``Block all public access'' and updated the \textbf{Bucket
        Policy} to allow public reads:

\begin{lstlisting}[style=jsonStyle]
{
  "Version": "2012-10-17",
  "Statement": [{
    "Effect": "Allow",
    "Principal": "*",
    "Action": "s3:GetObject",
    "Resource": "arn:aws:s3:::amzn-s3-cs5525-tonyn-bucket/*"
  }]
}
\end{lstlisting}

  \item Uploaded \texttt{index.html} to the bucket via the S3 console
        \textbf{Upload} button.
  \item Verified the live site at the S3 website endpoint URL.
\end{enumerate}

\subsubsection{Screenshots}

\screenshot{aws/01_aws_console_login.png}{AWS Console --- logged in with account name visible}
\screenshot{aws/02_s3_bucket_static_hosting.png}{S3 bucket created and Static Website Hosting enabled}
\screenshot{aws/03_bucket_policy.png}{Bucket Policy configured for public access}
\screenshot{aws/04_file_upload.png}{\texttt{index.html} uploaded and listed in Objects tab}
\screenshot{aws/05_live_site.png}{Live site accessible at S3 endpoint}

% ─────────────────────────────────────────────────────────────────────────────
\subsection{2.\ Google Cloud Platform (30\%)}
% ─────────────────────────────────────────────────────────────────────────────

\textbf{Method:} Compute Engine VM with Apache2 HTTP Server \textit{(IaaS)}

\subsubsection{Steps Taken}

\begin{enumerate}
  \item Signed into the \textbf{Google Cloud Console} at
        \texttt{console.cloud.google.com}.
  \item Created a new project: \texttt{cs5525-tonyn}.
  \item Navigated to \textbf{Compute Engine} $\to$ \textbf{VM Instances} $\to$
        Created a new VM:
        \begin{itemize}
          \item \textbf{Machine type:} \texttt{e2-micro} (free-tier eligible)
          \item \textbf{Boot disk:} Ubuntu 22.04 LTS
          \item \textbf{Firewall:} Checked \textit{Allow HTTP traffic} and
                \textit{Allow HTTPS traffic}
        \end{itemize}
  \item SSH'd into the VM using the \textbf{Cloud Shell SSH} button.
  \item Installed and started the Apache2 web server:

\begin{lstlisting}[style=bashStyle]
sudo apt update && sudo apt upgrade -y
sudo apt install apache2 -y
sudo systemctl start apache2
sudo systemctl enable apache2
\end{lstlisting}

  \item Verified the Apache default page via the external IP.
  \item Transferred the \texttt{index.html} content to the web root:

\begin{lstlisting}[style=bashStyle]
sudo nano /var/www/html/index.html
# Pasted HTML content, saved with Ctrl+O, exited with Ctrl+X
\end{lstlisting}

  \item Confirmed the live page at \href{http://34.66.74.246}{\texttt{http://34.66.74.246}}.
\end{enumerate}

\subsubsection{Screenshots}

\screenshot{gcp/01_gcp_console_login.png}{GCP Console --- logged in with account email visible}
\screenshot{gcp/02_vm_instance_running.png}{Compute Engine VM instance running (external IP shown)}
\screenshot{gcp/03_ssh_apache_install.png}{SSH terminal --- Apache2 installation and systemctl output}
\screenshot{gcp/04_file_transfer_webroot.png}{SSH terminal --- writing \texttt{index.html} to \texttt{/var/www/html/}}
\screenshot{gcp/05_live_site.png}{Live site accessible at GCP external IP}

% ─────────────────────────────────────────────────────────────────────────────
\subsection{3.\ Microsoft Azure (30\%)}
% ─────────────────────────────────────────────────────────────────────────────

\textbf{Method:} Azure Blob Storage Static Website \textit{(PaaS)}

\subsubsection{Steps Taken}

\begin{enumerate}
  \item Signed into the \textbf{Azure Portal} at \texttt{portal.azure.com}.
  \item Created a \textbf{Storage Account}: \texttt{cs5525tonynstorage}
        \begin{itemize}
          \item \textbf{Region:} East US
          \item \textbf{Performance:} Standard
          \item \textbf{Redundancy:} LRS (Locally Redundant Storage)
        \end{itemize}
  \item Navigated to \textbf{Data management} $\to$ \textbf{Static website}.
  \item Enabled static website hosting and configured:
        \begin{itemize}
          \item \textbf{Index document name:} \texttt{index.html}
          \item \textbf{Error document path:} \texttt{404.html}
        \end{itemize}
  \item Noted the auto-generated primary endpoint URL.
  \item Navigated to the auto-created \textbf{\texttt{\$web}} container.
  \item Uploaded \texttt{index.html} via the \textbf{Upload} button in the
        container view.
  \item Verified the live site at the Azure static website endpoint.
\end{enumerate}

\subsubsection{Screenshots}

\screenshot{azure/01_azure_portal_login.png}{Azure Portal --- logged in with account name visible}
\screenshot{azure/02_storage_account_created.png}{Storage account \texttt{cs5525tonynstorage} created}
\screenshot{azure/03_static_website_enabled.png}{Static website feature enabled with index document set}
\screenshot{azure/04_file_upload_web_container.png}{\texttt{index.html} uploaded to the \texttt{\$web} container}
\screenshot{azure/05_live_site.png}{Live site accessible at Azure endpoint}

% =============================================================================
\section{Platform Comparison}
% =============================================================================

\begin{center}
\small
\begin{tabularx}{\linewidth}{@{} l X X X @{}}
\toprule
\textbf{Feature} & \textbf{AWS S3} & \textbf{GCP Compute Engine} & \textbf{Azure Blob Storage} \\
\midrule
Deployment Type     & PaaS (Managed)       & IaaS (VM)              & PaaS (Managed) \\
Setup Complexity    & Low                  & High                   & Low \\
Server Management   & None required        & Full (Apache2)         & None required \\
HTTPS Support       & Via CloudFront CDN   & Manual cert install    & Built-in \\
Free Tier           & 5 GB / 20K req       & 1$\times$ e2-micro/mo  & 5 GB / month \\
Deployment Speed    & $\sim$5 min          & $\sim$20 min           & $\sim$5 min \\
File Transfer       & S3 Console / CLI     & SCP / SSH / Shell      & Portal / Storage Explorer \\
IAM / Access        & Bucket Policy (JSON) & Firewall + IAM Roles   & Blob access + RBAC \\
Custom Domain       & Via Route 53         & Via DNS A record       & Via Azure CDN \\
\bottomrule
\end{tabularx}
\end{center}

% =============================================================================
\section{Epilog}
% =============================================================================

Deploying a simple personal website across AWS, Google Cloud Platform (GCP), and
Microsoft Azure was an eye-opening journey revealing both the similarities and
differences among those three giant web hosting providers.

AWS S3 and Azure Blob Storage stood out for their simplicity in hosting static
websites. I was able to enable static website hosting with a few clicks and set
the appropriate bucket or container policies to make my site publicly accessible.
The process was straightforward, and I had a live URL up and running in no time.
AWS S3, in particular, was enjoyable due to its simplicity and the clean,
easy-to-remember endpoint URL it provided.

GCP, on the other hand, took a more hands-on approach. I had to provision a
Compute Engine VM and manually install and configure Apache2. This method was
the most time-consuming but also the most educational. It gave me a deeper
understanding of how web servers operate at the infrastructure level. I learned
how to set up the server and how to configure the web server software.

The most challenging aspect across all three platforms was managing IAM policies,
bucket permissions, and firewall rules. A misconfigured firewall rule on GCP
initially blocked all HTTP traffic until I explicitly opened port 80 in the VPC
firewall settings. AWS's default ``Block Public Access'' setting required careful
policy overrides to serve content publicly. These experiences underscored the
importance of meticulous configuration management.

In my own reflection, I believe future classes should include a dedicated session
on cloud networking fundamentals, such as VPCs, security groups, and firewall
rules. A solid understanding of these concepts upfront would have significantly
reduced the troubleshooting time I spent across all three deployments. Further,
an initial exposure to infrastructure-as-code tools like Terraform could help
standardize and automate multi-cloud workflows and enhance both efficiency and
error reduction.

% =============================================================================
\section{Technologies Used}
% =============================================================================

\begin{itemize}
  \item \textbf{HTML / CSS} --- Static personal portfolio web page
  \item \textbf{Amazon S3} --- Object storage with static website hosting (AWS)
  \item \textbf{Google Compute Engine} --- Ubuntu 22.04 LTS virtual machine (GCP)
  \item \textbf{Apache2} --- HTTP web server installed on GCP VM
  \item \textbf{Azure Blob Storage} --- Object storage with static website hosting (Azure)
  \item \textbf{Google Cloud Shell / SSH} --- Remote VM access and file transfer
  \item \textbf{AWS Management Console} --- S3 bucket creation and management
  \item \textbf{Azure Portal} --- Storage account and container management
\end{itemize}

% =============================================================================
\section{Repository Structure}
% =============================================================================

\begin{lstlisting}[style=bashStyle, frame=none, backgroundcolor=\color{white}]
Assignment_2/
|-- index.html              # Web page deployed to all three platforms
|-- readme.md               # Project documentation
|-- assignment2.tex         # This LaTeX source file
|-- assignment2.pdf         # Compiled PDF report
`-- images/
    |-- aws/
    |   |-- 01_aws_console_login.png
    |   |-- 02_s3_bucket_static_hosting.png
    |   |-- 03_bucket_policy.png
    |   |-- 04_file_upload.png
    |   `-- 05_live_site.png
    |-- gcp/
    |   |-- 01_gcp_console_login.png
    |   |-- 02_vm_instance_running.png
    |   |-- 03_ssh_apache_install.png
    |   |-- 04_file_transfer_webroot.png
    |   `-- 05_live_site.png
    `-- azure/
        |-- 01_azure_portal_login.png
        |-- 02_storage_account_created.png
        |-- 03_static_website_enabled.png
        |-- 04_file_upload_web_container.png
        `-- 05_live_site.png
\end{lstlisting}

% =============================================================================
\section{Author}
% =============================================================================

\begin{center}
\begin{tabular}{rl}
  \textbf{Name:}    & Tony Nguyen \\
  \textbf{Degree:}  & Master of Science in Data Science and Analytics \\
  \textbf{University:} & University of Missouri -- Kansas City \\
  \textbf{Course:}  & CS 5525 --- Cloud Computing, Spring 2026 \\
\end{tabular}
\end{center}

\end{document}
